\section{議論と限界}

\paragraph{一次データソース側のデータ品質が測定の前提になる．}
当初はFortinet PSIRTのHTML変換データ\cite{vuls-data-raw-fortinet}を基準に測定していたが，このデータはCVRF/CSAFの更新に追従できておらず，2018〜2023年公開分ではノードの51\%でバージョンが空であった（実際のアドバイザリには記載がある）．
本稿の測定はすべてCVRF/CSAFに基づく．
一次データソースとの比較では，一次データソース側の抽出が健全であること自体を検証しなければならない．

\paragraph{製品の突合は綴りに依存する．}
CPEの製品部から区切り文字を除いて突合しているため，略称の差（\texttt{FortiClient Enterprise Management Server} と \texttt{forticlientems} 等）は自動では吸収できない．
表\ref{tab:match}のC2のうち「NVDに別のFortinet製品のCPEはある」ケース（Fortinet採番で98〜99\%）には，この対応付けの失敗が含まれる可能性がある．

\paragraph{観測期間の制約．}
4章は2026年7月29日時点の1点観測である．
5.5節のC4はGit履歴から測ったが，履歴はNVDが2025年6月26日，VulnCheckが2025年8月6日以降しか存在せず，それ以前の公開分は測れない．
また2026年後半の公開分は観測期間が短く，遅延の長い側が打ち切られる．

\paragraph{CPE適用条件の有無は検知可能性と等価ではない．}
4.3節および表\ref{tab:cna}のMicrosoftの例のとおり，条件が存在してもバージョンを特定できない場合がある．
逆にCPE辞書に製品が登録されていない場合は条件があっても検知できない．
本稿の付与率は検知可能性の上限を与えるものと解釈すべきである．

\paragraph{7章の「正解」は一次データソースである．}
未検知率・過剰検知率はDB-primaryを正解とみなした値であり，一次データソース側の過大な記載は過剰検知として現れる．
またMicrosoftのみ検知方式が異なり（適用済みKB対CPE），他の2ベンダと同列には比較できない．
スキャン対象は8環境と限られ，Microsoftの実機以外は宣言したCPEのみを持つ疑似サーバである．
またCiscoの製品名からCPEへの対応付けはアドバイザリの31.9\%しか覆っておらず，7章のCiscoの結果はIOSとIOS XEに限られる．

\paragraph{供給源の信頼度による抑制の代償と限界．}
VulnCheckのCPEは機械的に生成されるため，Microsoft Edgeのようにうまく機能するものと，Windowsのように大量の偽陽性を生むものが混在する（7.2節(2)(4)）．
Vulsが「検証済みソースが製品を定義していれば補助的ソースの検知を捨てる」規則を置いているのはこのためであり，規則そのものは必要な妥協である．
本稿が測ったのはその代償と限界である．
代償として，macOS 10.13.6では真陽性781件が捨てられている．
限界として，Windows 11 23H2の偽陽性833件はこの規則をすり抜けた．NVDが置いた抽象CPE（\texttt{all\_windows}）とVulnCheckが展開した具体CPE（\texttt{windows\_11\_23h2}）が別製品として扱われるためである．
規則を緩めれば前者は救えるが後者は悪化する．
必要なのは信頼度の調整ではなく，照合不能な抽象CPEをソース側で明示的に扱うことだと考えられる．

\paragraph{LLM利用に伴うリスク．}
6章のP1・P2は非決定性を生成時に閉じ込めるが，未実装のP3はデータ経路にLLMが残るため，キャッシュと検証器に決定性を依存する．
いずれの検証器も形式的な整合性を確認できるが「原文の意味を正しく解釈したか」は保証しない．
P1で出典間の相互照合により118件の不一致を検出したことは，単一の出力を信じてはならないことの裏返しでもある．
人手レビューのキューを設計に含めており，そのレビュー負荷が運用可能な水準に収まるかは未検証である．

その他の限界は付録\ref{app:limits}に記す．
